
\documentclass{article}
\usepackage{colortbl}
\usepackage{makecell}
\usepackage{multirow}
\usepackage{supertabular}

\begin{document}

\newcounter{utterance}

\centering \large Interaction Transcript for game `function\_detective', experiment `list\_passive\_examples', episode 91 with mistral{-}large{-}3{-}azure.
\vspace{24pt}

{ \footnotesize  \setcounter{utterance}{1}
\setlength{\tabcolsep}{0pt}
\begin{supertabular}{c@{$\;$}|p{.15\linewidth}@{}p{.15\linewidth}p{.15\linewidth}p{.15\linewidth}p{.15\linewidth}p{.15\linewidth}}
    \# & \multicolumn{2}{c}{Player} && \multicolumn{2}{c}{Game Master} \\
    \hline

    \theutterance \stepcounter{utterance}  
    & & & \multicolumn{4}{p{0.6\linewidth}}{
        \cellcolor[rgb]{0.9,0.9,0.9}{
            \makecell[{{p{\linewidth}}}]{
                \texttt{\tiny{[P1$\langle$GM]}}
                \texttt{I have a hidden Python function for you to figure out.} \\
\\ 
\texttt{Your job is to infer what it does and then submit a correct Python implementation.} \\
\\ 
\texttt{Function signature:} \\
\texttt{(items: List{[}int{]}) {-}> int} \\
\\ 
\texttt{You have 15 turns total.} \\
\\ 
\texttt{I control which valid input{-}output examples you see. You can request them one at a time.} \\
\\ 
\texttt{You must make your best SOLVE attempt no later than the second{-}to{-}last turn.} \\
\texttt{Do not wait until the very end.} \\
\\ 
\texttt{You can use these actions:} \\
\texttt{{-} JSON format described at the end of the prompt} \\
\texttt{{-} SOLVE: ```python ... ```} \\
\\ 
\texttt{Do not add import statements in your guess.} \\
\texttt{Anything you might need is already available, so just write the function body in hypothesis and return the answer.} \\
\\ 
\texttt{Preloaded examples:} \\
\texttt{1. ({[}5, {-}5, 5{]}, 50)} \\
\\ 
\texttt{Scoring:} \\
\texttt{{-} correctness is required for any non{-}zero score} \\
\texttt{{-} in interactive modes, solving with fewer information requests is better} \\
\\ 
\texttt{Input formatting reminders:} \\
\texttt{{-} strings must be quoted} \\
\texttt{{-} booleans must be True/False} \\
\texttt{{-} use commas between multiple values} \\
\texttt{{-} do not write any extra action tags in the same message} \\
\\ 
\texttt{Format example:} \\
\texttt{You write:} \\
\texttt{\{'current\_mode': <probing/confirming/don't know; choose whichever best describes your current state in the task>, 'input\_rationale': <brief description of why you chose this hypothesis>, 'hypothesis': <either the exact string "unknown" if you have no guess yet OR your current best guess as runnable Python matching the given signature, with a one{-}line docstring stating what it does. Never write placeholder code, if you cannot commit to a guess, use "unknown". Match the given signature.} \\
\texttt{e.g.} \\
\texttt{def f(x):} \\
\texttt{	"""Returns the absolute value of x."""} \\
\texttt{	return x if x >= 0 else {-}x} \\
\\ 
\texttt{Use "unknown" if you have no guess.>, 'confidence\_score': <score between 1{-}100. 100 means highest confidence in your hypothesis>\}} \\
\\ 
\texttt{Output only one valid json per turn.} \\
            }
        }
    }
     \\ \\

    \theutterance \stepcounter{utterance}  
    & \multicolumn{4}{p{0.6\linewidth}}{
        \cellcolor[rgb]{0.561,0.737,0.561}{
            \makecell[{{p{\linewidth}}}]{
                \texttt{\tiny{[P1$\rangle$GM]}}
                \texttt{```json} \\
\texttt{\{} \\
\texttt{    "current\_mode": "probing",} \\
\texttt{    "input\_rationale": "The preloaded example ({[}5, {-}5, 5{]}, 50) suggests the function might be calculating the product of the absolute values of the items multiplied together. For {[}5, {-}5, 5{]}, the absolute values are {[}5, 5, 5{]}, and their product is 125, but the output is 50, which doesn't match. Alternatively, it might be the sum of the squares of the items: 25 + 25 + 25 = 75, which also doesn't match. Another possibility is the product of the items multiplied by their count: (5 * {-}5 * 5) * 3 = {-}375, which doesn't match. The output 50 is 5 * 5 * 2, but this is unclear. I need another example to narrow it down.",} \\
\texttt{    "hypothesis": "unknown",} \\
\texttt{    "confidence\_score": 10} \\
\texttt{\}} \\
\texttt{```} \\
            }
        }
    }
    & & \\ \\

    \theutterance \stepcounter{utterance}  
    & & & \multicolumn{4}{p{0.6\linewidth}}{
        \cellcolor[rgb]{0.9,0.9,0.9}{
            \makecell[{{p{\linewidth}}}]{
                \texttt{\tiny{[P1$\langle$GM]}}
                \texttt{OUTPUT: ({[}{]}, 0)} \\
            }
        }
    }
     \\ \\

    \theutterance \stepcounter{utterance}  
    & \multicolumn{4}{p{0.6\linewidth}}{
        \cellcolor[rgb]{0.561,0.737,0.561}{
            \makecell[{{p{\linewidth}}}]{
                \texttt{\tiny{[P1$\rangle$GM]}}
                \texttt{```json} \\
\texttt{\{} \\
\texttt{    "current\_mode": "probing",} \\
\texttt{    "input\_rationale": "The new example ({[}{]}, 0) suggests the function handles empty lists by returning 0. The first example ({[}5, {-}5, 5{]}, 50) still doesn't fit simple patterns like sum of squares (75) or product of absolute values (125). Another possibility is the sum of the absolute values multiplied by the count of non{-}zero items: (5 + 5 + 5) * 3 = 45, which is close but not 50. Alternatively, it could be the sum of the absolute values of the items multiplied by the maximum absolute value in the list: (5 + 5 + 5) * 5 = 75, which doesn't match. I need another example to test hypotheses.",} \\
\texttt{    "hypothesis": "unknown",} \\
\texttt{    "confidence\_score": 20} \\
\texttt{\}} \\
\texttt{```} \\
            }
        }
    }
    & & \\ \\

    \theutterance \stepcounter{utterance}  
    & & & \multicolumn{4}{p{0.6\linewidth}}{
        \cellcolor[rgb]{0.9,0.9,0.9}{
            \makecell[{{p{\linewidth}}}]{
                \texttt{\tiny{[P1$\langle$GM]}}
                \texttt{OUTPUT: ({[}0{]}, 0)} \\
            }
        }
    }
     \\ \\

    \theutterance \stepcounter{utterance}  
    & \multicolumn{4}{p{0.6\linewidth}}{
        \cellcolor[rgb]{0.561,0.737,0.561}{
            \makecell[{{p{\linewidth}}}]{
                \texttt{\tiny{[P1$\rangle$GM]}}
                \texttt{```json} \\
\texttt{\{} \\
\texttt{    "current\_mode": "probing",} \\
\texttt{    "input\_rationale": "The new example ({[}0{]}, 0) confirms that the function returns 0 for lists containing zero or empty lists. The first example ({[}5, {-}5, 5{]}, 50) still doesn't fit simple patterns. Possible hypotheses now include:\textbackslash{}n1. Sum of absolute values multiplied by the count of non{-}zero elements: (5+5+5)*3=45 (doesn't match 50)\textbackslash{}n2. Sum of squares multiplied by the count: (25+25+25)*3=225 (doesn't match)\textbackslash{}n3. Product of non{-}zero elements multiplied by their count: (5*{-}5*5)*3={-}375 (doesn't match)\textbackslash{}n4. Sum of absolute values multiplied by the maximum absolute value: (5+5+5)*5=75 (doesn't match)\textbackslash{}n5. Sum of absolute values of pairwise products: |5*{-}5| + |5*5| + |{-}5*5| = 25+25+25=75 (doesn't match)\textbackslash{}n\textbackslash{}nI need another example to distinguish between these possibilities.",} \\
\texttt{    "hypothesis": "unknown",} \\
\texttt{    "confidence\_score": 30} \\
\texttt{\}} \\
\texttt{```} \\
            }
        }
    }
    & & \\ \\

    \theutterance \stepcounter{utterance}  
    & & & \multicolumn{4}{p{0.6\linewidth}}{
        \cellcolor[rgb]{0.9,0.9,0.9}{
            \makecell[{{p{\linewidth}}}]{
                \texttt{\tiny{[P1$\langle$GM]}}
                \texttt{OUTPUT: ({[}1, 1{]}, 2)} \\
            }
        }
    }
     \\ \\

    \theutterance \stepcounter{utterance}  
    & \multicolumn{4}{p{0.6\linewidth}}{
        \cellcolor[rgb]{0.561,0.737,0.561}{
            \makecell[{{p{\linewidth}}}]{
                \texttt{\tiny{[P1$\rangle$GM]}}
                \texttt{```json} \\
\texttt{\{} \\
\texttt{    "current\_mode": "probing",} \\
\texttt{    "input\_rationale": "The new example ({[}1, 1{]}, 2) suggests the function might simply return the sum of the absolute values of the items. For {[}1, 1{]}, this gives 2 (matches). For {[}5, {-}5, 5{]}, this would give 15 (doesn't match 50). Alternatively, it could be the sum of the items multiplied by their count: (1+1)*2=4 (doesn't match 2). Another possibility is the sum of the items squared: 1+1=2 (matches), but for {[}5, {-}5, 5{]} this would be 25+25+25=75 (doesn't match 50). The only consistent pattern so far is that empty lists and lists with 0 return 0. The first example still doesn't fit, so I need to test another hypothesis: perhaps the sum of absolute values multiplied by the count of non{-}zero elements, but {[}5, {-}5, 5{]} gives (5+5+5)*3=45 (doesn't match 50). I need one more example to confirm.",} \\
\texttt{    "hypothesis": "unknown",} \\
\texttt{    "confidence\_score": 40} \\
\texttt{\}} \\
\texttt{```} \\
\texttt{To help narrow it down, I'd like to request the output for `{[}2} \\
            }
        }
    }
    & & \\ \\

    \theutterance \stepcounter{utterance}  
    & & & \multicolumn{4}{p{0.6\linewidth}}{
        \cellcolor[rgb]{0.9,0.9,0.9}{
            \makecell[{{p{\linewidth}}}]{
                \texttt{\tiny{[P1$\langle$GM]}}
                \texttt{Game Violation: ParseError: Response is not valid JSON.} \\
\texttt{Output ONLY:} \\
\texttt{One JSON with confidence\_score, current\_mode, input\_rationale, and hypothesis keys ...} \\
\texttt{OR} \\
\texttt{SOLVE: ```python ... ```} \\
            }
        }
    }
     \\ \\

    \theutterance \stepcounter{utterance}  
    & \multicolumn{4}{p{0.6\linewidth}}{
        \cellcolor[rgb]{0.561,0.737,0.561}{
            \makecell[{{p{\linewidth}}}]{
                \texttt{\tiny{[P1$\rangle$GM]}}
                \texttt{```json} \\
\texttt{\{} \\
\texttt{    "current\_mode": "probing",} \\
\texttt{    "input\_rationale": "The examples so far suggest a pattern where the function might be calculating the sum of the absolute values of all items multiplied by the number of non{-}zero items. For {[}1, 1{]}, this gives (1+1)*2=4 (doesn't match 2). Alternatively, it could be the sum of the absolute values of the items multiplied by the maximum absolute value in the list. For {[}1, 1{]}, this gives (1+1)*1=2 (matches). For {[}5, {-}5, 5{]}, this would give (5+5+5)*5=75 (doesn't match 50). Another possibility is the sum of the squares of the items: {[}1,1{]} gives 1+1=2 (matches), but {[}5,{-}5,5{]} gives 25+25+25=75 (doesn't match 50). The only consistent pattern is that empty lists and lists with 0 return 0. I'll request one more example to confirm.",} \\
\texttt{    "hypothesis": "unknown",} \\
\texttt{    "confidence\_score": 50} \\
\texttt{\}} \\
\texttt{```} \\
            }
        }
    }
    & & \\ \\

    \theutterance \stepcounter{utterance}  
    & & & \multicolumn{4}{p{0.6\linewidth}}{
        \cellcolor[rgb]{0.9,0.9,0.9}{
            \makecell[{{p{\linewidth}}}]{
                \texttt{\tiny{[P1$\langle$GM]}}
                \texttt{OUTPUT: ({[}1, 2, 3{]}, 14)} \\
            }
        }
    }
     \\ \\

    \theutterance \stepcounter{utterance}  
    & \multicolumn{4}{p{0.6\linewidth}}{
        \cellcolor[rgb]{0.561,0.737,0.561}{
            \makecell[{{p{\linewidth}}}]{
                \texttt{\tiny{[P1$\rangle$GM]}}
                \texttt{```json} \\
\texttt{\{} \\
\texttt{    "current\_mode": "confirming",} \\
\texttt{    "input\_rationale": "The new example ({[}1, 2, 3{]}, 14) strongly suggests the function is calculating the sum of the squares of the items plus the sum of the products of all unique pairs. For {[}1, 2, 3{]}: (1² + 2² + 3²) + (1*2 + 1*3 + 2*3) = (1+4+9) + (2+3+6) = 14 (matches). For {[}5, {-}5, 5{]}: (25+25+25) + (5*{-}5 + 5*5 + {-}5*5) = 75 + ({-}25+25{-}25) = 75{-}25 = 50 (matches). For {[}1,1{]}: (1+1) + (1*1) = 2+1 = 3 (but output is 2), which doesn't fit. Alternative hypothesis: sum of squares of items plus sum of absolute values of pairwise products. For {[}1,2,3{]}: (1+4+9) + (2+3+6) = 25 (doesn't match). Another possibility: sum of squares of items multiplied by their count. For {[}1,2,3{]}: (1+4} \\
            }
        }
    }
    & & \\ \\

    \theutterance \stepcounter{utterance}  
    & & & \multicolumn{4}{p{0.6\linewidth}}{
        \cellcolor[rgb]{0.9,0.9,0.9}{
            \makecell[{{p{\linewidth}}}]{
                \texttt{\tiny{[P1$\langle$GM]}}
                \texttt{Game Violation: ParseError: Response is not valid JSON.} \\
\texttt{Output ONLY:} \\
\texttt{One JSON with confidence\_score, current\_mode, input\_rationale, and hypothesis keys ...} \\
\texttt{OR} \\
\texttt{SOLVE: ```python ... ```} \\
            }
        }
    }
     \\ \\

    \theutterance \stepcounter{utterance}  
    & \multicolumn{4}{p{0.6\linewidth}}{
        \cellcolor[rgb]{0.561,0.737,0.561}{
            \makecell[{{p{\linewidth}}}]{
                \texttt{\tiny{[P1$\rangle$GM]}}
                \texttt{```json} \\
\texttt{\{} \\
\texttt{    "current\_mode": "confirming",} \\
\texttt{    "input\_rationale": "The pattern now clearly points to the function calculating the sum of all pairwise products (including squares) in the list. For {[}1,2,3{]}: 1*1 + 1*2 + 1*3 + 2*1 + 2*2 + 2*3 + 3*1 + 3*2 + 3*3 = 1+2+3+2+4+6+3+6+9 = 36 (doesn't match 14). However, the sum of squares plus sum of unique pairwise products works for {[}1,2,3{]} and {[}5,{-}5,5{]}, but fails for {[}1,1{]}. The correct pattern appears to be: sum of squares of all elements plus sum of products of all unique pairs. For {[}1,1{]}: 1+1 + 1*1 = 3 (but output is 2). The only consistent pattern is: sum(x\_i * x\_j for all i <= j). For {[}1,1{]}: 1*1 + 1*1 + 1*1 = 3 (still doesn't match). Final hypothesis: sum of squares of elements plus sum of products of adjacent pairs (with wrap{-}around). For {[}1,1{]}: 1+1 + 1*1 = 3. This isn} \\
            }
        }
    }
    & & \\ \\

    \theutterance \stepcounter{utterance}  
    & & & \multicolumn{4}{p{0.6\linewidth}}{
        \cellcolor[rgb]{0.9,0.9,0.9}{
            \makecell[{{p{\linewidth}}}]{
                \texttt{\tiny{[P1$\langle$GM]}}
                \texttt{Game Violation: ParseError: Response is not valid JSON.} \\
\texttt{Output ONLY:} \\
\texttt{One JSON with confidence\_score, current\_mode, input\_rationale, and hypothesis keys ...} \\
\texttt{OR} \\
\texttt{SOLVE: ```python ... ```} \\
            }
        }
    }
     \\ \\

    \theutterance \stepcounter{utterance}  
    & \multicolumn{4}{p{0.6\linewidth}}{
        \cellcolor[rgb]{0.561,0.737,0.561}{
            \makecell[{{p{\linewidth}}}]{
                \texttt{\tiny{[P1$\rangle$GM]}}
                \texttt{```json} \\
\texttt{\{} \\
\texttt{    "current\_mode": "confirming",} \\
\texttt{    "input\_rationale": "After analyzing all examples, the function appears to calculate the sum of the squares of all elements plus the sum of the products of all unique pairs (i < j). For {[}5, {-}5, 5{]}: (25+25+25) + (5*{-}5 + 5*5 + {-}5*5) = 75 + ({-}25+25{-}25) = 50. For {[}1,2,3{]}: (1+4+9) + (1*2 + 1*3 + 2*3) = 14. For {[}1,1{]}: (1+1) + (1*1) = 3 (but output is 2), which suggests we should exclude the square terms. Final hypothesis: sum of products of all unique pairs (i < j). For {[}1,1{]}: 1*1 = 1 (doesn't match 2). The correct pattern is actually the sum of all possible products (including squares), which is equivalent to (sum(items))\^{}2. For {[}5,{-}5,5{]}: (5{-}5+5)\^{}2 = 25 (doesn't match 50). The actual solution is the sum of absolute values of all pairwise products (including squares). For {[}5,{-}5,5{]}: |2} \\
            }
        }
    }
    & & \\ \\

    \theutterance \stepcounter{utterance}  
    & & & \multicolumn{4}{p{0.6\linewidth}}{
        \cellcolor[rgb]{0.9,0.9,0.9}{
            \makecell[{{p{\linewidth}}}]{
                \texttt{\tiny{[P1$\langle$GM]}}
                \texttt{Game Violation: ParseError: Response is not valid JSON.} \\
\texttt{Output ONLY:} \\
\texttt{One JSON with confidence\_score, current\_mode, input\_rationale, and hypothesis keys ...} \\
\texttt{OR} \\
\texttt{SOLVE: ```python ... ```} \\
            }
        }
    }
     \\ \\

    \theutterance \stepcounter{utterance}  
    & \multicolumn{4}{p{0.6\linewidth}}{
        \cellcolor[rgb]{0.561,0.737,0.561}{
            \makecell[{{p{\linewidth}}}]{
                \texttt{\tiny{[P1$\rangle$GM]}}
                \texttt{```json} \\
\texttt{\{} \\
\texttt{    "current\_mode": "confirming",} \\
\texttt{    "input\_rationale": "After carefully analyzing all examples, the function appears to calculate the sum of the squares of all elements. However, this doesn't match {[}5,{-}5,5{]} (25+25+25=75≠50). The correct pattern is actually the sum of absolute values of all pairwise products (including each element with itself). For {[}5,{-}5,5{]}: |5*5| + |5*{-}5| + |5*5| + |{-}5*5| + |{-}5*{-}5| + |5*5| = 25+25+25+25+25+25=150 (doesn't match 50). The breakthrough came when realizing it's the square of the sum of absolute values: for {[}5,{-}5,5{]}, (5+5+5)\^{}2=225 (doesn't match). The actual solution is the sum of absolute values multiplied by the count of elements: (5+5+5)*3=45 (close to 50). Final realization: it's the sum of absolute values multiplied by the maximum absolute value in the list. For {[}5,{-}5,5{]}: (5+5+5)*5=75 (doesn't match). The correct answer is the sum of squares of absolute} \\
            }
        }
    }
    & & \\ \\

    \theutterance \stepcounter{utterance}  
    & & & \multicolumn{4}{p{0.6\linewidth}}{
        \cellcolor[rgb]{0.9,0.9,0.9}{
            \makecell[{{p{\linewidth}}}]{
                \texttt{\tiny{[P1$\langle$GM]}}
                \texttt{Game Violation: ParseError: Response is not valid JSON.} \\
\texttt{Output ONLY:} \\
\texttt{One JSON with confidence\_score, current\_mode, input\_rationale, and hypothesis keys ...} \\
\texttt{OR} \\
\texttt{SOLVE: ```python ... ```} \\
            }
        }
    }
     \\ \\

    \theutterance \stepcounter{utterance}  
    & \multicolumn{4}{p{0.6\linewidth}}{
        \cellcolor[rgb]{0.561,0.737,0.561}{
            \makecell[{{p{\linewidth}}}]{
                \texttt{\tiny{[P1$\rangle$GM]}}
                \texttt{SOLVE:} \\
\texttt{```python} \\
\texttt{def f(items):} \\
\texttt{    """Returns the sum of the squares of all elements plus the sum of the absolute values of all pairwise products."""} \\
\texttt{    total = 0} \\
\texttt{    for i in range(len(items)):} \\
\texttt{        for j in range(len(items)):} \\
\texttt{            total += abs(items{[}i{]} * items{[}j{]})} \\
\texttt{    return total // 2} \\
\texttt{```} \\
\\ 
\texttt{Wait, this doesn't match the examples. Let me try again with the correct solution:} \\
\\ 
\texttt{SOLVE:} \\
\texttt{```python} \\
\texttt{def f(items):} \\
\texttt{    """Returns half the square of the sum of absolute values of elements."""} \\
\texttt{    abs\_sum = sum(abs(x) for x in items)} \\
\texttt{    return (abs\_sum * abs\_sum) // 2} \\
\texttt{```} \\
\\ 
\texttt{No, this gives 45 for {[}5,{-}5,5{]} (15*15//2=112.5→112). Final attempt:} \\
\\ 
\texttt{SOLVE:} \\
\texttt{```python} \\
\texttt{def f(items):} \\
\texttt{    """Returns the sum of all pairwise products (including squares) of absolute values."""} \\
\texttt{    total = 0} \\
\texttt{    for i in range(len(items)):} \\
\texttt{        for j in range(i, len(items)):} \\
\texttt{            total += abs(items{[}i{]} * items{[}j{]})} \\
\texttt{    return total} \\
\texttt{```} \\
            }
        }
    }
    & & \\ \\

    \theutterance \stepcounter{utterance}  
    & & & \multicolumn{2}{p{0.3\linewidth}}{
        \cellcolor[rgb]{0.9,0.9,0.9}{
            \makecell[{{p{\linewidth}}}]{
                \texttt{\tiny{[GM$|$GM]}}
                \texttt{Game Verdict: LOSS} \\
            }
        }
    }
    & & \\ \\

    \theutterance \stepcounter{utterance}  
    & & & \multicolumn{2}{p{0.3\linewidth}}{
        \cellcolor[rgb]{0.9,0.9,0.9}{
            \makecell[{{p{\linewidth}}}]{
                \texttt{\tiny{[GM$|$GM]}}
                \texttt{Incorrect.} \\
\texttt{Failed Cases:} \\
\texttt{Input: {[}{[}5, {-}5, 5{]}{]} | Expected: 50 | Got: 112} \\
\texttt{Input: {[}{[}1, 2, 3{]}{]} | Expected: 14 | Got: 18} \\
\texttt{Input: {[}{[}{-}1, 0, 1{]}{]} | Expected: 1 | Got: 2} \\
            }
        }
    }
    & & \\ \\

    \theutterance \stepcounter{utterance}  
    & & & \multicolumn{2}{p{0.3\linewidth}}{
        \cellcolor[rgb]{0.9,0.9,0.9}{
            \makecell[{{p{\linewidth}}}]{
                \texttt{\tiny{[GM$|$GM]}}
                \texttt{def sum\_of\_positive\_squares(items: List{[}int{]}) {-}> int:} \\
\texttt{    return sum(x * x for x in items if x > 0)} \\
            }
        }
    }
    & & \\ \\

    \theutterance \stepcounter{utterance}  
    & & & \multicolumn{2}{p{0.3\linewidth}}{
        \cellcolor[rgb]{0.9,0.9,0.9}{
            \makecell[{{p{\linewidth}}}]{
                \texttt{\tiny{[GM$|$GM]}}
                \texttt{def f(items):} \\
\texttt{    """Returns the sum of the squares of all elements plus the sum of the absolute values of all pairwise products."""} \\
\texttt{    total = 0} \\
\texttt{    for i in range(len(items)):} \\
\texttt{        for j in range(len(items)):} \\
\texttt{            total += abs(items{[}i{]} * items{[}j{]})} \\
\texttt{    return total // 2} \\
            }
        }
    }
    & & \\ \\

    \theutterance \stepcounter{utterance}  
    & & & \multicolumn{2}{p{0.3\linewidth}}{
        \cellcolor[rgb]{0.9,0.9,0.9}{
            \makecell[{{p{\linewidth}}}]{
                \texttt{\tiny{[GM$|$GM]}}
                \texttt{Mode: passive\_examples} \\
\texttt{Accuracy: 0.050} \\
\texttt{Information requests: 4/15} \\
\texttt{Revealed passive examples: 5} \\
\texttt{Efficiency: 0.800} \\
\texttt{Slack: 11} \\
\texttt{Observed{-}pair consistency: 0.600} \\
\texttt{Parse errors: 5} \\
\texttt{Runtime errors: 0} \\
            }
        }
    }
    & & \\ \\

\end{supertabular}
}

\end{document}
